\documentclass[11pt,a4paper]{article}
\usepackage[T1]{fontenc}
\usepackage[utf8]{inputenc}
\usepackage{amsmath,amssymb,amsthm}
\usepackage[margin=1in]{geometry}
\usepackage[colorlinks=true,linkcolor=blue,urlcolor=blue]{hyperref}
\theoremstyle{plain}
\newtheorem{theorem}{Theorem}
\newtheorem{lemma}[theorem]{Lemma}
\newtheorem{proposition}[theorem]{Proposition}
\newtheorem{corollary}[theorem]{Corollary}
\theoremstyle{definition}
\newtheorem{remark}[theorem]{Remark}
% A quoted conjecture can be a single hundred-term formula whose terms are thirty-digit
% numbers. TeX cannot hyphenate those, so without a little stretch every such quote runs off
% the page; the linked-a-file papers are all of that shape.
\setlength{\emergencystretch}{3em}

\title{The empirical closed form for OEIS A220927, proved}
\author{Adrian Perez Fontelles\\ \small Independent researcher}
\date{1 September 2026}
\begin{document}
\maketitle

\begin{abstract}
OEIS A220927 carries an empirical closed form for $a(n)$ --- not a recurrence relating
consecutive terms but an explicit formula. It is true, and it is decidable rather than
empirical. The entry counts arrays subject to a condition on a bounded window of consecutive lines, so the lines of an array are the
vertices of a finite digraph, an array is a walk in it, and $a(n)$ is a walk count on
$S=391$ vertices. Any function of the form $\sum_j p_j(n)\lambda_j^{\,n}$ satisfies the
monic linear recurrence whose characteristic polynomial is
$\prod_j(x-\lambda_j)^{\deg p_j+1}$; here that polynomial is $q(x)=(x-1)^{39}$, of degree
$39$. Two things are then checked exactly: that the walk count satisfies the same
recurrence from some index on, which the Cayley--Hamilton theorem reduces to a finite
computation, and that the two sequences agree at $39$ consecutive indices past that
point. A monic recurrence determines every later term from its predecessors, so agreement
there is agreement for good.
\end{abstract}

\noindent\small 2020 Mathematics Subject Classification. 05A15, 11B37, 15A18.\normalsize

\section{The conjecture}

OEIS A220927 is ``Number of nX3 arrays of the minimum value of corresponding elements and their horizontal, vertical or antidiagonal neighbors in a random, but sorted with lexicographically nondecreasing rows and columns, 0..3 nX3 array.''. It has offset $1$ and begins
\[
10, 41, 283, 1448, 6687, 29847, 128985, 528548,\ \dots
\]
The entry states, as a conjecture contributed by its author and never marked settled:
\begin{quote}\raggedright\small
Empirical: a(n) = (1/523022617466601111760007224100074291200000000)*n\^{}38 + (1/9175835394150896697543986387720601600000000)*n\^{}37 - (1/43763920798811907937411699782451200000000)*n\^{}36 + (59/27555061243696386479111070233395200000000)*n\^{}35 + (307/1968218660264027605650790730956800000000)*n\^{}34 - (1549/57888784125412576636787962675200000000)*n\^{}33 + (1624193/994634563609361544032084085964800000000)*n\^{}32 + (121589/1961803872996768331424228966400000000)*n\^{}31 - (862211087/64169971845765260905295747481600000000)*n\^{}30 + (334940077/419411580691276215067292467200000000)*n\^{}29 + (881602781/386354510291963804027505868800000000)*n\^{}28 - (1123886873/366908366848968474859929600000000)*n\^{}27 + (15952682958133/81134447161312398845776232448000000)*n\^{}26 - (894295004698811/270448157204374662819254108160000000)*n\^{}25 - (228156596853291101/811344471613123988457762324480000000)*n\^{}24 + (5101733059113811/230561088835784026273873920000000)*n\^{}23 - (692780969776572349/1153803544130590192089169920000000)*n\^{}22 - (2243838262686431231/384601181376863397363056640000000)*n\^{}21 + (757680056062692907691431/731172092995226361139801620480000000)*n\^{}20 - (581959953416778427639843/16774528449283195869604085760000000)*n\^{}19 + (5511354224236467062364481117/17990681761856227570150381977600000000)*n\^{}18 + (40879216835452378507072451/2449711568880205279159910400000000)*n\^{}17 - (236826152119988143815376123/335377179072885246551654400000000)*n\^{}16 + (6642578071741399997758170841/619157869057634301326131200000000)*n\^{}15 + (28405985333085798680701919749417/760635442137303739179152179200000000)*n\^{}14 - (2151776489142623183168814040553/466075638564524349987225600000000)*n\^{}13 + (1424091974719336255283095572327437/17287169139484175890435276800000000)*n\^{}12 - (800612498869091029720619190616727/1920796571053797321159475200000000)*n\^{}11 - (4966763411298100637669295860759843/568398985311837982792089600000000)*n\^{}10 + (10551986050476453738601614474489679/55261012460873137215897600000000)*n\^{}9 - (3981523334369030919774590230531163843/2748084098835503719465574400000000)*n\^{}8 - (568854062299140169791514988575529/460547025110692763443200000000)*n\^{}7 + (190860909288639665111491234401381180257/1557247656006785441030492160000000)*n\^{}6 - (4867354810542169054981510999137717857/4325687933352181780640256000000)*n\^{}5 + (1778619289191249601556249307356327/363503187676653931146240000)*n\^{}4 - (253662472663677656114377414621/29473606296279365220000)*n\^{}3 - (178529187999606526995551083121/28268061177142554768000)*n\^{}2 + (1518270092662877319067/58075341924600)*n + 29455077 for n>15
\end{quote}
As of the ``Last modified'' line on the live entry (Oct 03 2025 19:07:39, revision 8) this is still
recorded as empirical, and nothing on the entry records it as proved. Write
\begin{multline*}
f(n)\;=\;\frac{n^{38}}{523022617466601111760007224100074291200000000} + \frac{n^{37}}{9175835394150896697543986387720601600000000} \\
- \frac{n^{36}}{43763920798811907937411699782451200000000} + \frac{59 n^{35}}{27555061243696386479111070233395200000000} \\
+ \frac{307 n^{34}}{1968218660264027605650790730956800000000} - \frac{1549 n^{33}}{57888784125412576636787962675200000000} \\
+ \frac{1624193 n^{32}}{994634563609361544032084085964800000000} + \frac{121589 n^{31}}{1961803872996768331424228966400000000} \\
- \frac{862211087 n^{30}}{64169971845765260905295747481600000000} + \frac{334940077 n^{29}}{419411580691276215067292467200000000} \\
+ \frac{881602781 n^{28}}{386354510291963804027505868800000000} - \frac{1123886873 n^{27}}{366908366848968474859929600000000} \\
+ \frac{15952682958133 n^{26}}{81134447161312398845776232448000000} - \frac{894295004698811 n^{25}}{270448157204374662819254108160000000} \\
- \frac{228156596853291101 n^{24}}{811344471613123988457762324480000000} + \frac{5101733059113811 n^{23}}{230561088835784026273873920000000} \\
- \frac{692780969776572349 n^{22}}{1153803544130590192089169920000000} - \frac{2243838262686431231 n^{21}}{384601181376863397363056640000000} \\
+ \frac{757680056062692907691431 n^{20}}{731172092995226361139801620480000000} - \frac{581959953416778427639843 n^{19}}{16774528449283195869604085760000000} \\
+ \frac{5511354224236467062364481117 n^{18}}{17990681761856227570150381977600000000} + \frac{40879216835452378507072451 n^{17}}{2449711568880205279159910400000000} \\
- \frac{236826152119988143815376123 n^{16}}{335377179072885246551654400000000} + \frac{6642578071741399997758170841 n^{15}}{619157869057634301326131200000000} \\
+ \frac{28405985333085798680701919749417 n^{14}}{760635442137303739179152179200000000} - \frac{2151776489142623183168814040553 n^{13}}{466075638564524349987225600000000} \\
+ \frac{1424091974719336255283095572327437 n^{12}}{17287169139484175890435276800000000} - \frac{800612498869091029720619190616727 n^{11}}{1920796571053797321159475200000000} \\
- \frac{4966763411298100637669295860759843 n^{10}}{568398985311837982792089600000000} + \frac{10551986050476453738601614474489679 n^{9}}{55261012460873137215897600000000} \\
- \frac{3981523334369030919774590230531163843 n^{8}}{2748084098835503719465574400000000} - \frac{568854062299140169791514988575529 n^{7}}{460547025110692763443200000000} \\
+ \frac{190860909288639665111491234401381180257 n^{6}}{1557247656006785441030492160000000} - \frac{4867354810542169054981510999137717857 n^{5}}{4325687933352181780640256000000} \\
+ \frac{1778619289191249601556249307356327 n^{4}}{363503187676653931146240000} - \frac{253662472663677656114377414621 n^{3}}{29473606296279365220000} \\
- \frac{178529187999606526995551083121 n^{2}}{28268061177142554768000} + \frac{1518270092662877319067 n}{58075341924600} \\
+ 29455077.
\end{multline*}

\section{The count is a walk count}

The entry's defining condition is local: it is a condition on a bounded window of consecutive lines. Take as vertices the
admissible configurations of such a window and put an edge from $u$ to $v$ when $v$ may
follow $u$, which the condition decides by inspection. An admissible array is then exactly a
walk, so with $M$ the adjacency matrix of that digraph on $S=391$ vertices, and $u,w$ the
vectors recording which windows may start and end an array,
\[
a(n)\;=\;u^{\!\top}M^{\,g(n)}w
\]
for the linear function $g$ that the entry's shape supplies. In particular $a$ is
$C$-finite: it satisfies some monic integer linear recurrence, though not necessarily a short
one.

\section{A closed form is a recurrence}

\begin{lemma}\label{lem:ann}
Let $f(m)=\sum_{j}p_j(m)\lambda_j^{\,m}$ with the $\lambda_j$ distinct and the $p_j$
polynomials, and put $q(x)=\prod_j(x-\lambda_j)^{\deg p_j+1}$. Then $q$ applied to the shift
operator annihilates $f$: writing $q(x)=x^{r}-\sum_{i=1}^{r}c_ix^{r-i}$,
\[
f(m)\;=\;\sum_{i=1}^{r}c_i\,f(m-i)\qquad\text{for every }m.
\]
\end{lemma}

\begin{proof}
The shift operator $E$ acts on $m\mapsto\lambda^{m}p(m)$ as
$(E-\lambda)\bigl(\lambda^{m}p(m)\bigr)=\lambda^{m+1}\bigl(p(m+1)-p(m)\bigr)$, and
$p(m+1)-p(m)$ has degree one less than $p$. So $(E-\lambda)^{\deg p+1}$ kills
$\lambda^{m}p(m)$, and the factors of $q$ commute.
\end{proof}

Here $q(x)=(x-1)^{39}$, of degree $39$; the closed form is a polynomial in $n$, so this is $(x-1)^{39}$, and the recurrence it encodes is
\begin{equation}\label{eq:rec}
a(n)\;=\;39\,a(n-1) - 741\,a(n-2) + 9139\,a(n-3) - 82251\,a(n-4) + 575757\,a(n-5) - 3262623\,a(n-6) + 15380937\,a(n-7) - 61523748\,a(n-8) + 211915132\,a(n-9) - 635745396\,a(n-10) + 1676056044\,a(n-11) - 3910797436\,a(n-12) + 8122425444\,a(n-13) - 15084504396\,a(n-14) + 25140840660\,a(n-15) - 37711260990\,a(n-16) + 51021117810\,a(n-17) - 62359143990\,a(n-18) + 68923264410\,a(n-19) - 68923264410\,a(n-20) + 62359143990\,a(n-21) - 51021117810\,a(n-22) + 37711260990\,a(n-23) - 25140840660\,a(n-24) + 15084504396\,a(n-25) - 8122425444\,a(n-26) + 3910797436\,a(n-27) - 1676056044\,a(n-28) + 635745396\,a(n-29) - 211915132\,a(n-30) + 61523748\,a(n-31) - 15380937\,a(n-32) + 3262623\,a(n-33) - 575757\,a(n-34) + 82251\,a(n-35) - 9139\,a(n-36) + 741\,a(n-37) - 39\,a(n-38) + a(n-39).
\end{equation}

\section{The proof}

\begin{lemma}\label{lem:crit}
With $q$ as above, $a$ satisfies \eqref{eq:rec} for all large $n$ if and only if
$u^{\!\top}M^{\,j}q(M)w=0$ for every $j\ge0$, and it is enough to check $j<S$.
\end{lemma}

\begin{proof}
Substituting the walk expression, the residual $a(n)-\sum_ic_ia(n-i)$ equals
$u^{\!\top}M^{\,g(n)-r}q(M)w$ up to the fixed linear reparametrisation $g$. For the bound,
$M^{S}$ is an integer combination of $I,M,\dots,M^{S-1}$ by Cayley--Hamilton, so the values
$u^{\!\top}M^{j}q(M)w$ satisfy the monic recurrence given by the characteristic polynomial of
$M$; $S$ consecutive zeros therefore force all later ones, and the last nonzero value pins the
threshold exactly.
\end{proof}

\begin{theorem}
$a(n)=f(n)$ for every $n\ge16$.
\end{theorem}

\begin{proof}
The vector $w'=q(M)w$ was formed by $39$ matrix--vector products in exact integer
arithmetic, and $u^{\!\top}M^{j}w'$ evaluated for $j=0,1,\dots$, again exactly. Those integers
vanish from the point corresponding to $n=55$ onwards, and the run was continued until the
working vector was identically zero, which settles every larger $j$ at once. So $a$ satisfies
\eqref{eq:rec} for every $n>54$. By Lemma~\ref{lem:ann} $f$ satisfies \eqref{eq:rec}
identically. Both were then evaluated in exact arithmetic at the $39$ consecutive indices
$n=55,\dots,93$ and agree at each. A monic recurrence of order $39$
determines $a(m)$ from $a(m-1),\dots,a(m-39)$, and for every $m\ge94$ both
$m>54$ and $m-39\ge55$ hold, so induction on $m$ gives $a(m)=f(m)$ for every
$m\ge55$.
Below $n=55$ the recurrence argument gives nothing, since \eqref{eq:rec} is only
established for $a$ from $n=55$ on. The remaining indices $n=16,\dots,54$ were
therefore checked one at a time, directly against the walk count in exact integer arithmetic,
and agree.
\end{proof}

The range is tight in the sense that was checked: at $n=15$ the closed form and the sequence differ, so no larger range is available.

\section{Verification}

Four checks, each independent of the algebra above.

First, the walk model reproduces every one of the $22$ terms the entry publishes,
exactly, in integer arithmetic. This is what ties the model to the entry: the model is built
by reading the entry's English, and a misreading gives a different digraph and different
counts.

Second, the closed form was evaluated in exact rational arithmetic --- not floating point ---
at every published term from $n=16$ on, and agrees with the entry's own DATA at each.

Third, the annihilation test of Lemma~\ref{lem:crit} was carried out exactly, and the model
and the closed form were then compared directly at sixty consecutive indices beyond
$n=16$, far past where the proof already settles the matter.

Fourth, the closed form was deliberately perturbed --- by a constant and by a term linear in
$n$ --- and each perturbed form was rejected by the same comparison, so the test is not one
that anything would pass.

\begin{thebibliography}{9}
\bibitem{oeis} The OEIS Foundation, \emph{The On-Line Encyclopedia of Integer
Sequences}, \url{https://oeis.org}, 2026. Sequence A220927.
\bibitem{stanley} R.~P.~Stanley, \emph{Enumerative Combinatorics}, Volume 1, 2nd ed.,
Cambridge University Press, 2011. (Transfer-matrix method, Section 4.7; $C$-finite sequences,
Section 4.1.)
\bibitem{kp} M.~Kauers and P.~Paule, \emph{The Concrete Tetrahedron}, Springer, 2011.
(Closed forms and linear recurrences with constant coefficients.)
\bibitem{horn} R.~A.~Horn and C.~R.~Johnson, \emph{Matrix Analysis}, 2nd ed., Cambridge
University Press, 2013.
\end{thebibliography}

\end{document}
